\section{FOL Semantics}

For the semantics of FOL we will want to leverage as much as possible from our
treatment of TFL semantics. But FOL sentences are not atomic symbols of truth
values, they are constructed from basic FOL expressions using the rules of
grammar. We will need to find a way to account for the truth of atomic sentences
which rests on the meanings of these basic parts.

We will use \emph{models} to give a semantics to first-order logic. They will
be responsible for associating the expressions of first-order logic with a
domain of individuals.

\begin{definition}[Jobs for Models]
  \label{def:jobs-for-model}
  An \emph{model} for first-order logic will:
  \begin{enumerate}
  \item determine the \emph{truth} of sentences
  \item set a non-empty \emph{domain};
  \item assign each name a \emph{reference} in the domain; and
  \item assign each predicate an \emph{extension} made from items in the domain.
  \end{enumerate}
\end{definition}

Before we start formalizing models, let's look at some of the tools we will
need. First, a notational convention. We will use \elib\ealpha\erib to talk
about the reference (or extension) of an expression \ealpha. This will be useful
in definitions and metalogic, where the semantic objects are often central to
the discussions and a full description of them would be quite long-winded
indeed.

\subsection{The Basic Idea of Sets}

We will start by informally introducing a new formal mechanism, sets. For our
purposes, a \emph{set} is a collection of objects. The domain is a set. We can
symbolize a set of anything (symbols, objects, etc) using curly braces: ${a,b}$.
We will not dig more deeply into the properties and behaviors of sets, other
than to note that they have no order or repeated items. When all of the elements
of one set \eDelta and found in another set \eGamma, we say that \eDelta is a
\emph{subset} of \eGamma ($\eDelta\esubset\eGamma$). When some item $x$ is
found in a set \eGamma, we will say that $x$ is an \emph{element} of \eGamma
($x\in\eGamma$).\footnote{Often when we use $x\in\eGamma$, we mean to refer
to $x$ not to the relation $x\in\eGamma$, per se. We just want to note that
we are talking about an $x$ from \eGamma.}

Sets are a technical device that constrain the models of FOL. Like truth values
in TFL, the extensions of predicates (i.e. sets) are the closest we can get to
their \emph{meaning}. Beyond its subject belonging to the set of all sunsets,
all the nuances of meaning in a predicate like ``is a sunset'' are lost. Where
in TFL two sentences that have the same truth value in all situations are
equivalent, in FOL two predicates that have the same extension are equivalent
too.

\subsection{Domains, Reference and Extension}

First, let's consider the meanings of the various basic expressions of FOL. We
can skip brackets, since their meaning straightforwardly relates the rules of
grammar and the scope of operators.

\begin{definition}[Domain]
  \label{def:fol-domain-set}
  The \emph{domain} of an FOL model is the set of all objects relevant to the
  model. A domain for first-order logic must include an individual for each of
  the names in the language, but multiple names may refer to the same
  individual.
\end{definition}

We can define a domain either \emph{extensionally}, by listing all of its
individuals, or with a \emph{domain description}.

The object that a term picks out is called its \emph{reference}.

\begin{definition}[Reference]
\label{def:fol-reference}
We take singular terms to \emph{refer} to individuals in the domain. We will
call the individual that a name refers to its \emph{reference}. Although
many names may pick out the same reference, each name only picks out one.
\end{definition}

Two terms can refer to the same object.

\begin{definition}[Coreference]
  \label{def:fol-coreference}
  Two terms $\etau_1$ and $\etau_2$ are \emph{co-referential} iff they have the
  same reference.
\end{definition}

Likewise, we interpret a predicate as its \emph{extension}.

\begin{definition}[Extension of a Predicate]
  \label{def:fol-extension}
  The \emph{extension} of a predicate in first-order logic is the set
  of objects to which that predicate applies. If \ebeta is a 1-place predicate,
  then the extension of \ebeta is a subset of the domain. If \ebeta is an
  n-place relation, then its extension is a set of n-place \emph{tuples} of
  items from the domain
\end{definition}

We equate predicates on the basis of their extension.

\begin{definition}[Coextension]
  \label{def:fol-coextension}
  Two predicates $\ebeta_1$ and $\ebeta_2$ are  \emph{co-extensional} iff they
  have the same extension.
\end{definition}

We won't linger on sets and tuples, just note that they offer convenient
notational devices. One way we could represent an FOL interpretation is
extensionally: by listing the contents of the domain and the extensions of all
the predicates by rote. Tuples make it possible for us to do that for n-place
predicates where n>1. 

\begin{grayarg}
\item[\textbf{D:}] \{ John, Mary, Fred, Susie \}
  \item[$j$:] John
  \item[$m$:] Mary
  \item[$f$:] Fred
  \item[$s$:] Susie
  \item[$P$:] \{Fred,Susie\}
  \item[$E$:] \{Fred,Susie\}
  \item[$L$:] \{\textlangle Fred,Mary\textrangle,\textlangle Mary,Fred\textrangle,\textlangle Susie,John\textrangle,\textlangle John,John\textrangle\}
\end{grayarg}

We get a quick sense of who is a philosopher ($P$), who studies epistemology
($E$), and who loves whom ($L$). But this FOL symbolization does not let us
capture any of the differences in meaning between \emph{philosopher} and
\emph{specialist in epistemology}. This has particularly notable consequences
when dealing with so-called \emph{empty predicates}: those predicates that do
not have any objects in their extensions. On our approach, they all have the
same meaning.

\subsection{Setting Reference in Descriptive Domains}

Now consider this descriptively defined interpretation which is less explicit
about the contents of the domain and the extension of the predicate.

\begin{argument}
  \item[\textbf{D:}] all of the birds of prey in Chemung County
  \item[$h$:] Hank
  \item[$T$:] lives at Tanglewood
\end{argument}

While the domain refers to a \textbf{specific} set of individuals, we don't
exactly have a name for every bird in the county. This makes the explicit
definition of our predicates much more challenging. How do we talk about all
of the birds in the domain who don't have a name? Although there are some
tricks for doing so in the metalogic, the short version is that we can't,
especially when we start talking about infinite sets.

Things get especially weird when we try to figure out if two objects are
\emph{distinct}. The classic example is:

\begin{argument}
  \item[$p$:] Phosphorus
  \item[$h$:] Hesperus
  \item[$v$:] Venus
\end{argument}

All of these are names for the same thing. We will stick to this version of
\emph{direct reference} in our understanding of FOL semantics.\footnote{

  In future logic courses, especially \emph{Modal Logic}, we will see how terms
  that refer to the same thing have funny behavior in so-called \emph{epistemic
  contexts}.
  
  \begin{argument}
  \item[ ] The ancients knew that Phosphorus rose in the morning.
  \item[ ] Phospherus is Hesperus.
  \item[\therefore] The ancients knew that Hesperus rose in the morning.
  \end{argument}
  
  Although substitution of equivalents is typically truth-preserving. This is
  not a valid inference, but formalizing that proves to be difficult.
  Ultimately, all three names refer to one item. Our domain should reflect that.
  But how do we know when two items in the domain are the distinct? One
  proposal is that they are never the same. Specifically, each item is just
  itself. If two things are in the domain, then they really are two different
  things. We will wait until our survey of metalogic in Logic 2 to expand on
  this point. Suffice it to say, we will assume that the substitution of
  expressions with the same reference (or extension) will be truth-preserving
  in FOL, just as it is in TFL. But, technically, this is only correct in an
  \emph{extensional context} (which is often circularly defined as a context
  where truth-preservation for the substitution of equivalents is a correct
  assumption).
  
}

One trick is to assume that the elements of the domain can, in principle, be
counted (and in a consistent order). This gives us a way to make unique
reference to each of the items in a set: the natural numbers. As in
\cite{Mag25} we will use numbers to stand for objects in the domain. Names are
taken to point \emph{directly} at the objects these numbers represent, and in
a \emph{rigid} way: a name always refers to \textbf{the same} object.

Complicating matters further, we cannot assume that domains are finite. There
are a number of relations that cannot be validated by anything other than an
infinite set of individuals. For example, when a \emph{connected}\footnote{Every individual is related to some other individual.} two-place predicate $R$ is also both \emph{transitive}\footnote{Rxy and Ryz always means Rxz.} and \emph{irreflexive}\footnote{No individual is related to itself.}, this is only possible in an infinite domain. In other words, the following three wffs cannot be all valid together in a finite model.

\begin{center}
  \begin{tabular}{c c}
    connected    &  \eforall x\eexists yRxy              \\
    transitive   &  \eforall x((Rxy \eand Ryz) \eif Rxz) \\
    irrelflexive &  \eforall x\enot Rxx                  \\
  \end{tabular}
\end{center}

This certainly complicates any potential search for counter-examples. We will
need more powerful validation tools then those appropriate to finite models.
Ultimately, we will need to reason about the models themselves.

\subsection{Interpretation in First-Order Logic}

\begin{definition}[FOL Model]
  We define a model for first-order logic as comprising \textbf{D} and
  $\elib\cdot\erib$ where:
  \begin{outline}[enumerate]
  \1 \textbf{D} is a non-empty \emph{domain of individuals}; and
  \1 $\elib\cdot\erib$ is an \emph{interpretation function} which assigns each
  term a reference and each predicate an extension as follows:
  \2 \elib\etau\erib is a \emph{reference} in \textbf{D} for each term \etau;
  \2 \elib\ephi\erib is a \emph{truth-value} for each 0-place predicate \ephi;
  \2 \elib\ebeta\erib is a \emph{set of individuals} from \textbf{D} for each
  1-place predicate \ebeta;
  \2 \elib\ebeta\erib is a \emph{set of n-tuples of individuals} from \textbf{D}
  for each n-place predicate \ebeta, when n>1
  \end{outline}
\end{definition}

We've already seen how we can explicitly specify these things with sets of
objects (and tuples of objects). Using numbers for our domain, we can recast
the above example as follows.

\begin{grayarg}
  \item[\textbf{D:}] \{1,2,3,4\}
  \item[$j$:] 1 (John)
  \item[$m$:] 2 (Mary)
  \item[$f$:] 3 (Fred)
  \item[$s$:] 4 (Susie)
  \item[$P$:] \{3,4\}
  \item[$E$:] \{3,4\}
  \item[$L$:] \{\textlangle 3,2\textrangle,\textlangle 2,3\textrangle,\textlangle 4,1\textrangle,\textlangle 1,1\textrangle\}
\end{grayarg}

The sets are easy enough to read, but the tuple notation for relations is
sometimes difficult to scan. Instead, we can diagrammatically represent our
interpretation of the relation $L$.

\begin{figure}[h]
  \centering
  \begin{tikzpicture}[
      node distance=2cm,
      every node/.style={minimum size=0.5cm, inner sep=2pt}, % circle, draw
      ->,
      >=Stealth
    ]
    \node (1) {1};
    \node (2) [right of=1] {2};
    \node (3) [below of=2] {3};
    \node (4) [left of=3] {4};
    \draw (3) -- (2);
    \draw (2) -- (3);
    \draw (4) -- (1);
    \draw (1) to [loop left, looseness=6] (1);
  \end{tikzpicture}
  \caption{An interpretation of $L$}
  \label{fig:model-of-love}
\end{figure}

\subsection{Assessing Truth with the Resources of FOL}

Now that we have a handle on interpretations, we can start to talk about how
truth works in first-order logic. Recall from our informal survey that we need
to account for three types of sentences.

\begin{enumerate}
\item \textbf{atomic sentences:} sentences built from FOL terms and predicates
  but not from quantifers or connectives
\item \textbf{truth-functional sentences:} sentences whose main operator is
  truth-functional
\item \textbf{quantified sentences:} sentences whose main operator is
  a quantifier of FOL
\end{enumerate}

Following the convention from \cite{Mag25}, when a given object is found in the
extension of a predicate, we will say that the predicate is \emph{true of} that
object. Similarly, we will say that a relation is \emph{true of} a tuple of
objects whenever that tuple occurs in the extension of the relation.

\subsubsection{Modelling Quantified Sentences}

We'll adopt the model from \cite{Mag25}[Chapter 31] to keep things simple. But
we will adopt the notation from \cite{Smi20}, where we label slots with numbers
like \textcircled{1}. It is important to note that \textcircled{1} and 1 are
different symbols of the metalanguage. \textcircled{1} is a placeholder for
an empty term slot in a predicate. It will be filled by a term or a variable
by the grammar. 1 is just the symbol we've decided to use to stand for a
particular object in our domain, specifically Aristotle.

\begin{argument}
\item[\textbf{D:}] everyone born before 2000 CE
\item[a:] 1 (Aristotle)
\item[b:] 2 (Beyonce)
\item[P]: \textcircled{1} is a philosopher
\item[R:] \textcircled{1} was born before \textcircled{2}
\end{argument}

Note that we are not actually giving our names and predicates extensional
definitions here. We will assume truth based on real-world conditions. When
we construct our truth definition for quantifier-governed sentences, we will
want our rules to give some specific results on this interpretation, namely:

\begin{argument}
\item[ ] ``Someone was born before Beyonce'' is true
\item[ ] ``Some philosopher was born before Aristotle'' is true
\item[ ] ``Everyone is born before someone'' is false
\end{argument}

symbolized as:

\begin{argument}
\item[ ] $\eexists xR(b,x)$ is true
\item[ ] $\eexists x(P(x)\eand R(x,a))$ is true
\item[ ] $\eforall x\eexists y(R(x,y))$ is false
\end{argument}

Let's see how we can get these results, starting with some basic rules for the
truth of atomic sentences.

\subsubsection{Atomic Sentences}

We can now offer the following definitions of truth in first-order logic. As
noted above, we will depart slightly from \cite{Mag25} and use the notation
$\elib\ealpha\erib$ to mean the reference (or extension) of the FOL expression
\ealpha, given a model.

The sentence $\ebeta\etau_1,\dots,\etau_n$ is true in an model $\mathcal{M}$
just in case the references of $\ebeta\etau_1,\dots,\etau_n$ form an n-tuple
that can be found in the extension of \ebeta.

\begin{definition}[Truth for Atomic Sentences]
  \label{def:fol-atomic-truth}
  Let $\mathcal(M)=\langle\boldsymbol{D},\elib\cdot\erib\rangle$ be a model for
  FOL.

  $\elib\ebeta\etau_1,\dots,\etau_n\erib_\mathcal(M) = T$
  iff $\langle\elib\etau_1\erib_\mathcal{M},\dots,
  \elib\tau_n\erib_\mathcal{M}\rangle\in \elib\beta\erib_\mathcal{M}$
  \footnote{From now on, we will drop the $\mathcal{M}$ subscript whenever the
  choice of model is clear from context.}
\end{definition}

The sentence $\etau_1=\etau_2$ is true in an intepretation iff $\|\etau_1\|_I$
and $\|\etau_2\|_I$ are the very same individual in the domain $n\ein D$.

\begin{definition}[Truth for Identity Sentences]
  \label{def:fol-identity-truth}
  $\elib\etau_1=\etau_2\erib = T$ iff $\elib\etau_1\erib=\elib\etau_2\erib$  
\end{definition}

\subsubsection{Truth-Functional Sentences}

For truth-functional sentences, we can just import our connective definitions
from TFL.

\begin{definition}[Truth for Connective-Governed Sentences]
  \label{def:fol-connective-truth}
  The sentential connectives in FOL work just as they do in TFL:\\
  \begin{enumerate}
  \item $\elib\enot\ephi\erib=T$ iff $\elib\ephi\erib=F$ 
  \item $\elib\ephi\eand\epsi\erib=T$ iff $\elib\ephi\erib=T$ and
    $\elib\epsi\erib=T$    
  \item $\elib\ephi\eand\epsi\erib=T$ iff either $\elib\ephi\erib=T$ or
    $\elib\epsi\erib=T$ (or both)
  \item $\elib\ephi\eif\epsi\erib=T$ iff either $\elib\ephi\erib=F$
    or $\elib\epsi\erib=T$
  \item $\elib\ephi\eiff\epsi\erib=T$ iff $\elib\ephi\erib=\elib\epsi\erib$
  \end{enumerate}  
\end{definition}

\subsubsection{Quantified Sentences}

The definition of truth for atomic sentences (\ref{def:fol-atomic-truth}) puts
a lot of responsibility on the interpretation function and term expressions. We
need to be able to pick out the references for all of our terms in order to
check the extensions of the predicates and determine the truth of atomic
sentences. But what about the cases where we filled in term slots with something
other than a term, like a variable? How do we evaluate the truth of
(well-formed) formulas of this type? This is where the dummy names and dummy
sentences come in.

Before defining truth for quantified sentences, we need to pause once
more and consider the role that names play in definition
\ref{def:fol-atomic-truth}. Specifically, we are assuming that each term
picks out a specific item in the domain. For names this is straightfoward (by
direct reference, the interpretation of a name just \textbf{is} its object). But
what about variables? Here, we can leverage the manner in which we built our
domain (i.e. as \emph{countable}, and therefore ``named'' by the natural
numbers).

How does this work in practice? One soulution is to use the psuedo-names
(numbers) associated with each domain item to talk about possible changes we
might make to the interpretation. Specifically, we note that we \textbf{could}
extend the interpretation by adding a new name and being precise about its
reference. Indeed, one assumes that we do this sort of thing in ordinary English
all the time (when we learn the name of a person we've seen before, for
example).

In order to define our rules for truth, we will lean pretty heavily on the
observation above that two equivalent expressions (i.e. two expressions with
the same reference or extension) can be substituted for each other in a FOL
sentence without changing the truth value of the sentence.

\subsubsection{Interlude: Semantics with Open Sentences}

As much as possible we want our definition to work at the sentence level, but
when we are dealing with free variables, we have to acommodate formulas other
than sentences as well.

Recall from Section \ref{sec:fol-free-variables} that we have chosen to use a
slightly different grammar than \cite{Mag25}. Specifically, our grammar will
not create any \emph{open sentences}: sentences with one or more \emph{unbound
variables} (\emph{free variables}). The approach in \cite{Mag25} builds
quantified sentences from open sentences by adding a quantifier with a variable
that matches the free one. They get their open sentences by allowing variables
to act like terms in the syntax: a variable can fill any slot in a predicate as
though it was a term.

The important point for the semantics is that these formulas (i.e. open
sentences) exist. They are built along the way as part of the process for
constructing quantified sentences. As such, the semantics in \cite{Mag25} can
talk about them as though they are well-defined. Indeed, they are.

To see how the semantics of quantified sentences works in \cite{Mag25} (which is
a common approach in modern predicate logic), we will define open sentences as
a type of \emph{subformula}. From that definition, we can explore the standard
treatment before we look at the adjustments needed for our official definitions.

\begin{definition}[Open Sentence]
  In our first-order logic, only terms and predicates can be used to construct
  atomic sentences. By definition, each variable is bound by the quantifier
  that introduces it. Neverthelss, we can think of open sentences as subformulas
  of quantified sentences.

  Let \etheta be a quantifier symbol (\eforall or \eexists) and \etheta\exi\ephi
  be a quantified sentence of FOL. We define \ephi(\exi) as a \emph{subformula}
  of \etheta\exi\ephi and further stupulate that it:
  \begin{enumerate}
  \item is an open sentence: it contains zero or more free occurrences of the
    variable \exi (but typically one or more);
  \item is not a well-formed formula in our grammar
  \item ...
  \end{enumerate}
\end{definition}

With open sentences at our disposal, we can now present the lightly-modified
definition of \emph{substitution instance} from \cite{Mag25}. In the traditional
definition of truth for quantified sentences, we use substitution instances to
talk about their \emph{satisfaction}.

\begin{definition}[Substitution Instance]
  \label{def:fol-substitution-instance}
  Let \ephi be a formula of first-order logic. Also let \exi be a variable
  and \etau be a term. We define \ephi(\exi) to mean that \ephi has zero or
  more free occurrences of \exi and \ephi(\etau) to be the result of
  replacing all of the free occurrences of \exi in \ephi with \etau. The
  resulting formula is called an \emph{substitution instance} of \ephi and
  \etau is called its \emph{instantiating name}.
\end{definition}

Armed with this notation, we can now talk explicitly about updating our
interpretations with exactly the names we need to assess the truth of our
quantified sentences. Let's start informally.

Recall that our model contains information about each of its names, specifically
which $n\in D$ we are talking about. We can imagine adding more entries of this
type by stipulating a name for any of the objects in \textbf{D}. We will call
this practice \emph{updating} an interpretation.

\begin{definition}[Interpretation Update]
  \label{def:fol-update}
  Let $\mathcal{M}=\langle\mathbf{D},\elib\cdot\erib\rangle$ be an FOL model,
  \edelta be a name and $n$ be an object in \textbf{D}. Then
  $\elib\cdot\erib^{[n/\edelta]}$ is just like $\elib\cdot\erib$ except 
  that it assigns the object $n$ to the name \edelta.
\end{definition}

Then we can talk about whether or not a given object $n$ is true of a predicate
$\ebeta$. We will call this unmediated truth of objects in these new sentences
\emph{satisfaction} and the new name we add a \emph{dummy name}.

\begin{definition}[Satisfaction]
  \label{def:fol-satisfaction}
  Let $\mathcal{M}=\langle\mathbf{D},\elib\cdot\erib\rangle$ be an FOL model.
  
  Then an object $n\in\mathbf{D}$ \emph{satisfies} \ephi(\exi) in $\mathcal{M}$
  iff $\elib\ephi(\edelta)\erib^{[n/\edelta]}=T$, where the \emph{dummy name}
  \edelta does not already occur in \ephi(\exi).

  We will refer to \ephi(\etau) in these cases as a \emph{dummy sentence}.
\end{definition}

More broadly, \ephi(\etau) is an \emph{instance} of \eexists\exi\ephi(\exi).

Now that we've done the metalogical heavy lifting in our definition of
satisfaction, our truth definitions for quantifier-governed sentences in
straight-forward.

\begin{definition}[Truth for Quantifier-Governed Sentences]
  \label{def:fol-quantifier-truth-mag}
  Let $\mathcal{M}=\langle\mathbf{D},\elib\cdot\erib\rangle$ be an FOL model.
  \begin{enumerate}
  \item $\elib\eforall\exi\ephi(\exi)\erib=T$ iff
    every $n\in\mathbf{D}$ satisfies $\ephi(\exi)$
  \item $\elib\eexists \exi\ephi(\exi)\erib=T$ iff at least one
    $n\in\mathbf{D}$ satisfies $\ephi(\exi)$
  \end{enumerate}
\end{definition}

\subsubsection{Quantified Truth without Open Sentences}

We can reframe the definitions in \ref{def:fol-quantifier-truth-mag} in terms
of updating the interpretation function. There is at least one way to update
$\elib\cdot\erib$ to make \eexists\exi\ephi true; and \eforall\exi\ephi is true
no matter how we update $\elib\cdot\erib$.

\begin{definition}[Quantifier Truth and Interpretation Updates]
  \label{def:fol-quantifier-truth}
  Let $\mathcal{M}=\langle\mathbf{D},\elib\cdot\erib\rangle$ be an FOL model.
  
  We define the truth of quantified sentences as follows:
  
  \begin{enumerate}
  \item $\elib\eforall\exi\ephi\erib=T$ iff
    $\elib\ephi\erib^{[n/\exi]}=T$ for every $n\in\mathbf{D}$
  \item $\elib\eexists \exi\ephi\erib=T$ iff
    $\elib\ephi\erib^{[n/\exi]}=T$ for some $n\in\mathbf{D}$
  \end{enumerate}
  
\end{definition}

Note how we can formulate this definition without any reference to open
sentences like \ephi(\exi). We can just talk directly about $\elib\cdot\erib$.
These are, of course, just two ways of saying the same thing. We will adopt this
second definition because it aligns better with our syntax and the semantics
of open sentences seems incoherent, at best.

%% TODO: can we move this to the syntax? its form is not semantic.
\subsection{Quantifier Duals}

One of the interesting properties of the Square of Opposition [ref?] is that it
arranges the quantified sentences into pairs, or \emph{duals}, which differed
only in terms of negation. ``Every'' is a dual of ``Not every''. It is also a
dual of ``Every not''. These two flavors of duality depend on two different
flavors of negation: \emph{external} and \emph{internal}, respectively. Among
the interesting details of this arrangement is that the existential and
universal quantifiers are doubly dual: they are the result of both internal and
external negation.

\begin{observation}[Duals of Quantified Sentences]
  \label{obs:quantifer-duals}
  The quantified sentences of FOL can be related via the following equivalences.
  \begin{enumerate}
    \item ``Some'' is synonymous with ``Not every not''
    \item ``Every'' is synonymous with ``Not some not''
  \end{enumerate}
\end{observation}

\subsection{Metasemantics of FOL}

We can summarize the truth semantics of FOL as follows.

\begin{definition}[Truth in First-Order Logic]
  \label{def:fol-truth}
  Let $\mathcal(M)=\langle\boldsymbol{D},\elib\cdot\erib\rangle$ be a model for
  FOL. We can assess the truth of FOL sentences as follows.
  \begin{outline}[enumerate]
    \1 $\elib\ebeta\etau_1,\dots,\etau_n\erib=T$ iff
    $\langle\elib\etau_1\erib,\dots,\elib\tau_n\erib\rangle\in \elib\beta\erib$
    \1 $\elib\etau_1=\etau_2\erib=T$ iff $\elib\etau_1\erib=\elib\etau_2\erib$
    \1 $\elib\enot\ephi\erib=T$ iff $\elib\ephi\erib=F$ 
    \1 $\elib\ephi\eand\epsi\erib=T$ iff $\elib\ephi\erib=T$ and
    $\elib\epsi\erib=T$
    \1 $\elib\ephi\eand\epsi\erib=T$ iff either $\elib\ephi\erib=T$ or
    $\elib\epsi\erib=T$ (or both)
    \1 $\elib\ephi\eif\epsi\erib=T$ iff either $\elib\ephi\erib=F$
    or $\elib\epsi\erib=T$
    \1 $\elib\ephi\eiff\epsi\erib=T$ iff $\elib\ephi\erib=\elib\epsi\erib$
    \1 $\elib\eforall\exi\ephi\erib=T$ iff
    $\elib\ephi\erib^{[n/\exi]}=T$ for every $n\in\mathbf{D}$
    \1 $\elib\eexists \exi\ephi\erib=T$ iff
    $\elib\ephi\erib^{[n/\exi]}=T$ for some $n\in\mathbf{D}$
  \end{outline}

\end{definition}

With these truth definitions in place, we can now talk about how to formalize
the semantic concepts relevant to first-order logic. As we have thus far, we
will try to preserve as much as possible from our truth-functional definitions.

A set of premises \eGamma \emph{entails} a sentence \epsi only if it is never
the case that \eGamma are all false, but \epsi is true. 

\begin{definition}[Entailment]
  \label{def:fol-entailment}
  The sentences $\ephi_i,\dots,\ephi_n$ \emph{entail} the sentence \epsi
  ($\ephi_i,\dots,\ephi_n\eentails\epsi$)
  iff there is no model $\mathcal{M}=\langle\mathbf{D},\elib\cdot\erib\rangle$
  where $\elib\ephi_i\erib=F,\dots,\elib\ephi_n\erib=F$ and $\elib\epsi\erib=T$.
\end{definition}

A sentence \ephi is \emph{valid} iff \ephi is true on every interpretation. We
can call such \ephi a \emph{validity}.

\begin{definition}[Validity]
  \label{def:fol-validity}
  A sentence \ephi is \emph{valid} iff $\elib\ephi\erib=T$ for every
  model $\mathcal{M}=\langle\mathbf{D},\elib\cdot\erib\rangle$.
  \\
\end{definition}

A sentence \ephi is a \emph{contradiction} iff \ephi is false on every
interpretation.

\begin{definition}[Contradiction]
  \label{def:fol-contradiction}
  A sentence \ephi is a \emph{contradiction} iff $\elib\ephi\erib=F$ for every
  model $\mathcal{M}=\langle\mathbf{D},\elib\cdot\erib\rangle$.
\end{definition}

Similarly, an argument is valid just in case there is no interpretation
in which all of its premises are true and its conclusion is false. (i.e.
$\ephi_i,\dots,\ephi_n\eentails\epsi$).

\begin{definition}[Valid Argument]
  \label{def:fol-valid}
  The argument $\ephi_i,\dots,\ephi_n\therefore\epsi$ is \emph{valid} in
  first-order logic iff there is no model
  $\mathcal{M}=\langle\mathbf{D},\elib\cdot\erib\rangle$ where
  $elib\ephi_i\erib=T,\dots,\elib\ephi_n\erib=T$ and $\elib\epsi\erib=F$.
  Otherwise it is \emph{invalid}.
\end{definition}

Two FOL sentences \ephi and \epsi are \emph{equivalent} when they are true in
exactly the same interpretations (i.e. when both $\ephi\eentails\epsi$ and
$\epsi\eentails\ephi$).

\begin{definition}[Equivalent]
  \label{def:fol-equivalent-sentences}
  Two FOL sentences \ephi and \epsi are \emph{equivalent} iff
  $\elib\ephi\erib=\elib\epsi\erib$ for all models $\mathcal{M}=
  \langle\mathbf{D},\elib\cdot\erib\rangle$
\end{definition}

The FOL sentences $\ephi_1,\dots,\ephi_n$ are \emph{jointly satisfiable} when
there is an interpretation that makes them all true.

\begin{definition}[Joint Satisfiability]
  \label{def:fol-joint-satisfiability}
  The FOL sentences $\ephi_1,\dots,\ephi_n$ are \emph{jointly satisfiable} iff
  $\elib\ephi_i\erib=T$ for all $1\le i\le n$ on some model
  $\mathcal{M}=\langle\mathbf{D},\elib\cdot\erib\rangle$. Otherwise they are
  \emph{jointly unsatisfiable}.
\end{definition}

\subsection{Interpretations and Proofs}

\subsubsection{Proof and Counter-Models}

Disproving the semantic properties discussed in the last section amounts to
finding an interpretation that conflicts with the conditions in the definition.
Those requirements in the are summarized in the following table.

\begin{center}
  \begin{tabular}{c c}
    \textbf{to disprove} & \textbf{find an interpretation where} \\
    \hline \\
    validity & \ephi is false \\
    contradiction & \ephi is true \\
    equivalent & \ephi and \epsi differ in truth value \\
  \end{tabular}
\end{center}

We will call such a refutory interpretation a \emph{counter-model}. Generalizing
to multi-sentence arguments we can provide the following set of interlocking
definitions.

\begin{definition}[Effects of Counter-Models on Arguments]
  \label{def:fol-countermodel-effects}
  Let $\mathcal{M}=\langle\mathbf{D},\elib\cdot\erib\rangle$ be an FOL model.

  If $\elib\ephi_1\erib=T,\dots,\elib\ephi_n\erib=T$ but $\elib\epsi\erib=F$,
  then:\\
  \begin{enumerate}
  \item $\ephi_1,\dots,\ephi_n\therefore\epsi$ is \emph{invalid}; and
  \item $\ephi_1,\dots,\ephi_n\enotentails\epsi$; and
  \item $\ephi_1,\dots,\ephi_n,\enot\epsi$ are \emph{jointly satisfiable}
  \end{enumerate}
\end{definition}

For the examples in the rest of this section, we will assume the following
model.

\begin{grayarg}
\item[\textbf{D:}] \{1, 2, 3\}
\item[s:] 1 (Socrates)
\item[p:] 2 (Plato)
\item[ ]  3 (Aristotle)
\item[F:] \{1,3\}
\item[G:] \{ \} 
\item[L:] \{\textlangle 1,2\textrangle,\textlangle 2,3\textrangle,
  \textlangle 2,1\textrangle,\textlangle 3,3\textrangle\}
\end{grayarg}

The same model can be used to show none of the following are valid.

\begin{argument}
\item[ ] Fs\therefore\eforall xFx
\item[ ] \eexists x(Fx \eor Gx)\therefore\eexists xGx
\item[ ] \eforall x\eexists yLxy\therefore\eexists y\eforall xLxy
\end{argument}

First, this model is a counter-model for the validity of $Fm\eif\eforall xFx$.
To see this we can use the tools introduced in the last two lectures. First, a
parse tree reveals the deeper structure, including the specific instances that
need to be evaluated for truth.

\begin{center}
  \begin{forest}
    for tree={ grow=south, l sep=12pt, s sep=14pt }
    [{$Fs\eif\eforall xFx$} [{Fs}] [{\eforall xFx} [{Fc}] ] ] 
  \end{forest}
\end{center}

We need to show that $Fs$ is true on a countermodel but that $Fc$ is not, for
some \elib c\erib\ein\textbf{D}. This is indeed the case since 2 (Plato) is not
$F$ in our model.

We can evaluate the second example with the following parse tree:

\begin{center}
  \begin{forest}
    for tree={ grow=south, l sep=12pt, s sep=14pt }
    [{\eexists x(Fx \eor Gx)\eif\eexists xGx}
      [{\eexists x(Fx \eor Gx)}[{$Fa_1 \eor Ga_1$}[{$Fa_1$}][{$Ga_1$}]]]
      [{\eexists xGx}[{$Ga_2$}]]
    ] 
  \end{forest}
\end{center}

Once again, the antecedant is true. We could assign either 1 or 3 to the name
$a_1$, since both satisfy $Fx$ and hence $Fx\eor Gx$. We just need to find out
if anybody is $G$. Nobody is. The sentence is not valid.

\begin{center}
  \begin{forest}
    for tree={ grow=south, l sep=12pt, s sep=14pt }
    [{$\eforall x\eexists yLxy\eif\eexists y\eforall xLxy$}
      [{$\eforall x\eexists yLxy$} [{$\eexists yLcy$}]]
      [{$\eexists y\eforall xLxy$} [{$\eforall xLxa$}]]
    ] 
  \end{forest}
\end{center}

To show the sentence to be invalid we need a true antecendent ($\eexists yLcy$)
and a false consequent ($\eforall xLxa$). A look at the interpretation of $L$
will help.

\begin{center}
\begin{tikzpicture}[
    node distance=2cm,
    every node/.style={minimum size=0.5cm, inner sep=2pt}, % circle, draw
    ->,
    >=Stealth
  ]
  \node (1) {1};
  \node (2) [right of=1] {2};
  \node (3) [below of=2] {3};
  \draw (1) -- (2);
  \draw (1) -- (3);
  \draw (2) -- (1);
  \draw (3) to [loop below, looseness=6] (3);
\end{tikzpicture}
\end{center}

The antecedent is true: every node has an arrow going out (i.e. everyone loves
someone or other). The consequent, however, is false. There is no one that is
loved by everyone. 3 is the only individual that loves themselves, and 1 does
not love them. This is a counter-model.
$\eforall x\eexists yLxy\eif\eexists y\eforall xLxy$ is not valid.

\subsubsection{Discovering Countermodels}

% working backward to a counter-model (pp350-1)

The countermodel in the last section was given to us in advnace of the review.
But how can we find one ourselves among the potentially infinite possibilities?
The answer is to keep it simple. Start from a \emph{minimal model}: a model with
just one individual and no predicates. Then add in details as needed. It's more
art than science, so let's look at a simple example. Later we will prove with
natural deduction that $\eexists y\eforall xLxy\eproves\eforall x\eexists yLxy$.
Now, consider this proof via counter-model that
$\eforall x\eexists yLxy\enotentails\eexists y\eforall xLxy$

If we start with just one individual (\textbf{D}=\{1\}) then we are left with
only reflexive $L$=\{\textlangle 1,1\textrangle\} or empty $L$=\{\}. In the
first case the conclusion is true so the entailment is valid. In the second case
the premise is false so the entailment is valid. This is not a counter-model.
But it we are careful we should be able to find a countermodel with two
individuals (\textbf{D}=\{1,2\}). We want both individuals to be loved, so as
to sasify the premise, so we will connect them to each other. We want each
individual to be unloved by someone to make the conclusion false, so we will
just make that the person themselves (i.e. no one loves themself).

The resulting interpretation for $L$ can be depicted as:

\begin{figure}[h]
  \label{mod:minimal-symmetric}  
  \begin{center}
    \begin{tikzpicture}[
        node distance=2cm,
        every node/.style={minimum size=0.5cm, inner sep=2pt}, % circle, draw
        ->,
        >=Stealth
      ]
      \node (1) {1};
      \node (2) [right of=1] {2};
      \draw (1) -- (2);
      \draw (2) -- (1);
    \end{tikzpicture}
  \end{center}
  \caption{A simple counter-model proving
    $\eforall x\eexists yLxy\enotentails\eexists y\eforall xLxy$}
\end{figure}
  
\subsection{Proving Validity}

But proving the opposite (i.e. that there is no counter-model) requires more
substantial reasoning about all possible models, since we need to know that
our condition persists no matter what interpretation we use. We will look more
closely at this kind of metalogical proof in \emph{Logic 2}.

%% TODO: some examples of metalogical reasoning about models (pp349-50)

Ultimately, we can summarize the proof requirements for the various semantic
properties of FOL in a table, just as we did for TFL.

\begin{center}
  \begin{tabular}{c|c c}
    property & \textbf{yes} & \textbf{no} \\
    \hline
    validity       & all models & one counter-model \\
    contradiction  & all models & one counter-model \\
    equivalent     & all models & one counter-model \\
    unsatisfiable  & all models & one counter-model \\
    valid argument & all models & one counter-model \\
    entailment     & all models & one counter-model \\
  \end{tabular}
\end{center}

\subsection{** Bonus Material: Properties of Relations **}

Finally, we can leverage regular properties of relations to characterize some
predicates without fully specifying their extensions. For instance we know
that the following argument is valid.

\begin{argument}
\item[ ] Andy is taller than Bill.
\item[ ] Bill is taller than Charlie.
\item[\therefore] Andy is taller than Charlie. 
\end{argument}

This is because the \emph{taller than} relation is transitive. In other words:

\begin{example}
  $eforall x\eforall y\eforall z[(R(x,y)\eand R(y,z))\eif R(x,z)]$
\end{example}

Note that we have defined our relation in the familiar language of first-order
logic. A handful of useful properties of relations are thus defined in the
following table.

\begin{center}
  \begin{tabular}{c c}
    property & fol definition \\
    \hline
    transitive      &
    $eforall x\eforall y\eforall z[(R(x,y)\eand R(y,z))\eif R(x,z)]$\\
    reflexive       &
    $\eforall x(R(x,x))$ \\
    symmetric       &
    $\eforall x\eforall y(R(x,y)\eif R(y,x))$ \\
    anti-symmetric  &
    $\eforall x\eforall y((R(x,y)\eand x\neq y)\eif\enot R(y,x))$ \\
  \end{tabular}
\end{center}

